\documentclass{article}
\usepackage{babel}
\selectlanguage{english}
\usepackage{amsmath}
\usepackage{amssymb}
\usepackage[colorlinks=true, linkcolor=blue]{hyperref}
\usepackage{xcolor}
\usepackage[top=2cm, bottom=2cm,
            left=1.5cm, right=1.5cm]{geometry}
\usepackage{listings}
\usepackage{url}
\usepackage{graphicx}
\usepackage{cite}

% --- Inline style for code-level names (AthenaK variables, files, flags) ------
\definecolor{varbg}{HTML}{EEF1F7}
\definecolor{varfg}{HTML}{273C75}
\newcommand{\var}[1]{%
  {\setlength{\fboxsep}{1.5pt}%
   \colorbox{varbg}{\textcolor{varfg}{\texttt{\small #1}}}}}

% --- Run-in headers for the individual steps of a workflow -------------------
% Small caps + rule instead of \textbf, so they read as a level below
% \subsubsection rather than as another heading of the same weight.
\definecolor{stepfg}{HTML}{1A6B62}
\newcommand{\step}[1]{%
  \par\medskip\noindent
  {\color{stepfg}\rule[-0.15ex]{2pt}{1.5ex}\hspace{0.5em}\scshape #1}%
  \hspace{0.75em}\ignorespaces}

% Define short-commands for the variables
\newcommand{\gxx}{\var{adm\_gxx}}
\newcommand{\gxy}{\var{adm\_gxy}}
\newcommand{\gxz}{\var{adm\_gxz}}
\newcommand{\gyy}{\var{adm\_gyy}}
\newcommand{\gyz}{\var{adm\_gyz}}
\newcommand{\gzz}{\var{adm\_gzz}}
\newcommand{\vx}{\var{velx}}
\newcommand{\vy}{\var{vely}}
\newcommand{\vz}{\var{velz}}
\newcommand{\alp}{\var{z4c\_alpha}}
\newcommand{\betx}{\var{z4c\_betax}}
\newcommand{\bety}{\var{z4c\_betay}}
\newcommand{\betz}{\var{z4c\_betaz}}
\newcommand{\Bx}{\var{bcc1}}
\newcommand{\By}{\var{bcc2}}
\newcommand{\Bz}{\var{bcc3}}
\newcommand{\mr}{\mathrm}

\begin{document}
\section*{AthenaK post-processing computations and conventions}
\textit{This document summarizes the computations and conventions needed
to post-process a set of variables/quantities. It also collects some of the output
variables that are given in the code and how they are defined.}

\subsection*{Conventions}
\step{Variables}
\begin{itemize}
  \item Metric: The spatial metric is defined with lower indices, i.e. $\gamma_{ij}$. It can be
  extracted from the code by specifying the output variables \gxx, \gxy, \gxz,
  \gyy, \gyz, and \gzz or by giving the block of variables \var{adm}, which
  also contains the extrinsic curvature $K_{ij}$. Note that the $\langle\mathrm{adm}\rangle$
  block needs to exist for this to apply.
  \item Lapse: As a scalar quantity, the lapse $\alpha$ follows no convention and can be fetched by specifying
  the \alp variable in the respective $\langle\mathrm{output}\rangle$ block. Note that the
  $\langle\mathrm{z4c}\rangle$ block needs to exist for this to apply.
  \item Shift: Together with the lapse $\alpha$, the shift completes the set of gauge variables and is defined
  with upper indices, i.e. $\beta^i$. It can be specified in the output via \betx, \bety,
  and \betz. Note that also here the $\langle\mathrm{z4c}\rangle$ block needs to exist.
  \item Magnetic field components: In AthenaK, we evolve the densitized magnetic field components, i.e.
  $\tilde{B}^i=\sqrt{\gamma}B^i$ with upper indices. By specifying the output for the magnetic field components
  via \Bx, \By, and \Bz, we therefore mean $\tilde{B}^i$ cell-centered on the grid. $\tilde{B}^i$ is
  defined in the Eulerian frame and needs to be undensitized in the post-processing step. This requires the metric (see above).
  Note that in this case the $\langle\mathrm{mhd}\rangle$ block needs to exist.
  \item 3-velocity: The convention for the 3-velocity in AthenaK is a bit tricky. When we specify \vx, \vy, and
  \vz in the output, then we mean the Lorentz-weighted 3-velocity in the fluid frame $\tilde{u}^i=Wv^i$, where
  $v^i$ is the 3-velocity in the Eulerian frame and $W$ the Lorentz factor. In GR-Athena++ the variable $\tilde{u}^i$ gets
  directly evolved and is defined as the projection of the 4-velocity $u^{\mu}$ onto the hypersurface $\Sigma_t$. This is important
  to know when computing the 4-velocity in some post-processing steps below. For the output of the Lorentz-weighted 3-velocity
  either $\langle\mathrm{hydro}\rangle$ or $\langle\mathrm{mhd}\rangle$ needs to exist.
  \item M1 variables: The neutrino energy density $E$ follows no convention as a scalar quantity and can be extracted via
  \var{rad\_m1\_E}. On the other hand, the neutrino number flux $F_i$ is given with lower indices as \var{rad\_m1\_F}, and lastly
  the neutrino number density also follows no convention as \var{rad\_m1\_N}. All the aforementioned variables are given per species
  in the Eulerian frame. To obtain all of them in the fluid frame, one would have to project the stress-energy tensor accordingly.
\end{itemize}

\step{Derived variables}
Here we collect a couple of derived variables, that are needed for the computation of some post-processing
quantities below:
\begin{itemize}
  \item Lorentz-factor: The Lorentz factor in terms of the AthenaK conventions can be computed via
  \begin{align}
    W = \sqrt{1 + \tilde{u}^i\tilde{u}_i} = \sqrt{1 + \gamma_{ij}\tilde{u}^i\tilde{u}^j}\ |\ \tilde{u}^i = Wv^i, \nonumber
  \end{align}
  which is equal to the common form by realizing that
  \begin{align}
    W = \sqrt{1 + \tilde{u}^i\tilde{u}_i} = \sqrt{1 + W^2v^2} = \sqrt{1 + \frac{v^2}{1-v^2}} = \sqrt{\frac{1-v^2+v^2}{1-v^2}} = \sqrt{\frac{1}{1-v^2}}.\nonumber
  \end{align}
  Writing this in terms of the AthenaK variables gives
  \begin{align}\label{Equ:lorentz}
    \gamma_{ij}\tilde{u}^i\tilde{u}^j =& \gxx\vx\vx+2\gxy\vx\vy+2\gxz\vx\vz\\
    &+\gyy\vy\vy+2\gyz\vy\vz+\gzz\vz\vz,\nonumber\\
    W =& \sqrt{1+\gamma_{ij}\tilde{u}^i\tilde{u}^j} \nonumber.
  \end{align}
  \item Determinant of the spatial metric: To get the determinant of the spatial metric $\gamma_{ij}$, we can first write it
  in component form as
  \begin{align}
    \gamma = \gamma_{xx}\gamma_{yy}\gamma_{zz}-2\gamma_{xy}\gamma_{xz}\gamma_{yz}-\gamma_{xx}\gamma_{yz}\gamma_{yz}
                -\gamma_{yy}\gamma_{xz}\gamma_{xz}-\gamma_{zz}\gamma_{xy}\gamma_{xy}\nonumber .
  \end{align}
  In terms of the AthenaK variables this becomes
  \begin{align}\label{Equ:detgamma}
    \gamma =& \gxx\gyy\gzz-2\gxy\gxz\gyz-\gxx\gyz\gyz\\
    &-\gyy\gxz\gxz-\gzz\gxy\gxy.\nonumber
  \end{align}
  \item Inverse spatial metric components: Since AthenaK only provides $\gamma_{ij}$, we need to calculate the inverse spatial metric
  $\gamma^{ij}$ with help of the determinant of the metric which we can obtain via Equ.~\eqref{Equ:detgamma}. The components are then computed
  via
  \begin{align}
    \gamma^{xx}&=(\gamma_{yy}\gamma_{zz}-\gamma_{yz}^2)/\gamma,\nonumber\\
    \gamma^{yy}&=(\gamma_{xx}\gamma_{zz}-\gamma_{xz}^2)/\gamma,\nonumber\\
    \gamma^{zz}&=(\gamma_{xx}\gamma_{yy}-\gamma_{xy}^2)/\gamma,\nonumber\\
    \gamma^{xy}&=(\gamma_{xz}\gamma_{yz}-\gamma_{xy}\gamma_{zz})/\gamma,\nonumber\\
    \gamma^{xz}&=(\gamma_{xy}\gamma_{yz}-\gamma_{xz}\gamma_{yy})/\gamma,\nonumber\\
    \gamma^{yz}&=(\gamma_{xy}\gamma_{xz}-\gamma_{yz}\gamma_{xx})/\gamma.\nonumber
  \end{align}
  Or in terms of AthenaK variables
  \begin{align}
    \gamma^{xx}&=(\gyy\gzz-\gyz^2)/\gamma,\label{Equ:invgxx}\\
    \gamma^{yy}&=(\gxx\gzz-\gxz^2)/\gamma,\label{Equ:invgyy}\\
    \gamma^{zz}&=(\gxx\gyy-\gxy^2)/\gamma,\label{Equ:invgzz}\\
    \gamma^{xy}&=(\gxz\gyz-\gxy\gzz)/\gamma,\label{Equ:invgxy}\\
    \gamma^{xz}&=(\gxy\gyz-\gxz\gyy)/\gamma,\label{Equ:invgxz}\\
    \gamma^{yz}&=(\gxy\gxz-\gyz\gxx)/\gamma.\label{Equ:invgyz}
  \end{align}
  \item Undensitized Eulerian magnetic field components: Since AthenaK stores the densitized magnetic field
  components $\tilde{B}^i$, we need to undensitize them using Equ.~\eqref{Equ:detgamma}, i.e.
  \begin{align}
    B^i = \frac{\tilde{B}^i}{\sqrt{\gamma}}.\nonumber
  \end{align}
  Or in terms of AthenaK variables
  \begin{align}
    B^x &= \frac{\Bx}{\sqrt{\gamma}},\label{Equ:undensmag1}\\
    B^y &= \frac{\By}{\sqrt{\gamma}},\label{Equ:undensmag2}\\
    B^z &= \frac{\Bz}{\sqrt{\gamma}}\label{Equ:undensmag3}.
  \end{align}
  \item Magnetic field amplitude in the Eulerian frame: Using the metric components $\gamma_{ij}$ and the
  undensitized magnetic field components of Equs.~\eqref{Equ:undensmag1}-\eqref{Equ:undensmag3}, we can compute the
  magnetic field amplitude in the Eulerian frame via
  \begin{align}
    B^2 = \gamma_{ij}B^iB^j.\nonumber
  \end{align}
  Or in terms of the AthenaK variables
  \begin{align}\label{Equ:Bmag}
    B^2 =& \gxx\Bx\Bx+2\gxy\Bx\By+2\gxz\Bx\Bx\nonumber\\
    &+\gyy\By\By+2\gyz\By\Bz+\gzz\Bz\Bz.
  \end{align}
  \item Magnetic field amplitude in the fluid frame: To get the magnetic field amplitude in the fluid frame, we
  require again the metric components $\gamma_{ij}$, the undensitized magnetic field components $B^i$, and the
  Lorentz-weighted 3-velocity in the Lorentz frame $\tilde{u}^i$. In general, the magnetic field amplitude in the fluid
  frame can be written as
  \begin{align}\label{Equ:bmag}
    b^2 = \frac{(B^2 + (B^iWv_i)(B^iWv_i))}{W^2}.
  \end{align}
  We can make use of Equ.~\eqref{Equ:lorentz} and Equ.~\eqref{Equ:Bmag}. The remaining two terms can be computed
  as
  \begin{align}
    B^iWv_i = B^i\tilde{u}_i= B^i\gamma_{ij}\tilde{u}^j=&\gamma_{xx}B^x\tilde{u}^x+\gamma_{yy}B^y\tilde{u}^y+\gamma_{zz}B^z\tilde{u}^z\nonumber\\
              &+\gamma_{xy}(B^x\tilde{u}^y+B^y\tilde{u}^x)+\gamma_{xz}(B^x\tilde{u}^z+B^z\tilde{u}^x)+\gamma_{yz}(B^y\tilde{u}^z+B^z\tilde{u}^y)\nonumber .
  \end{align}
  Or in terms of the AthenaK variables:
  \begin{align}\label{Equ:BWv}
    B^i\gamma_{ij}\tilde{u}^i=&\gxx\vx\Bx+\gyy\By\vy+\gzz\Bz\vz \nonumber\\
    &+\gxy(\Bx\vy+\By\vx)+\gxz(\Bx\vz+\Bz\vx)\nonumber\\
    &+\gyz(\By\vz+\Bz\vy).
  \end{align}
  Combining Equs.~\eqref{Equ:lorentz},~\eqref{Equ:Bmag}, and~\eqref{Equ:BWv} in the definition for $b^2$, yields
  the desired result.
  \item 4-velocity with upper indices: To obtain the 4-velocity $u^{\mu}$, we require the metric components $\gamma_{ij}$, the 3-velocity $\tilde{u}^i$,
  and the gauge variables $\alpha$ and $\beta^i$. The definitions are
  \begin{align}
    u^t&=\frac{W}{\alpha},\nonumber\\
    u^i&=\tilde{u}^i-u^t\beta^i.\nonumber
  \end{align}
  By using Equs.~\eqref{Equ:lorentz}, we can use the AthenaK variables to write
  \begin{align}
    u^t&=\frac{W}{\alp},\label{Equ:vel4t}\\
    u^x&=\vx - u^t\betx,\label{Equ:vel4x}\\
    u^y&=\vy - u^t\bety,\label{Equ:vel4y}\\
    u^z&=\vz - u^t\betz\label{Equ:vel4z}.
  \end{align}
  \item 4-velocity with lower indices: To compute the 4-velocity with lower indices, we need the gauge variables
  $\alpha$ and $\beta^i$, as well as the metric components $\gamma_{ij}$ and 3-velocity. We start by writing down the general
  lowering operation
  \begin{align}
    u_t &= g_{t\nu}u^{\nu}=g_{tt}u^t+g_{ti}u^i=(-\alpha^2+\beta^k\beta_k)u^t+\beta_iu^i=-\alpha W+\beta^iu_i,\nonumber\\
    u_i &= g_{i\nu}u^{\nu}=g_{it}u^t+g_{ij}u^j=\beta_iu^t+\gamma_{ij}u^j=\gamma_{ij}\beta^ju^t+\gamma_{ij}u^j=\gamma_{ij}\left(u^j+\beta^ju^t\right).\nonumber
  \end{align}
  To reach the first equality, we have contracted with $\beta^i$ and realized from above that $u^t=\frac{W}{\alpha}$, i.e.
  \begin{align}
    u_t =-\alpha W+\beta^k\beta_ku^t+\beta_iu^i &= -\alpha W+\beta^i(\beta_iu^t)+\beta^i(\gamma_{ij}u^j)\nonumber\\
    &=\beta^i(\beta_iu^t+\gamma_{ij}u^j)\nonumber\\
    &=\beta^iu_i.\nonumber
  \end{align}
  In the second equality, we have relabeled and reordered the shift components in the first term, while in the second
  we have raised the index of the shift vector and also reordered. The last equality uses the general expression for $u_i$ above.
  We can also realize that the bracket in $u_i$, i.e. $u^j+\beta^ju^t$, is equal to the 3-velocity AthenaK uses, redefined in terms
  of the 4-velocity, see Equs.~\eqref{Equ:vel4x}-\eqref{Equ:vel4z}. Therefore, we have $u^j+\beta^ju^t=\tilde{u}^j$.
  We can therefore express the 4-velocity with lowered indices in terms of the AthenaK variables as
  \begin{align}
    u_x &= \gxx\vx+\gxy\vy+\gxz\vz,\label{Equ:vel4x_d}\\
    u_y &= \gxy\vx+\gyy\vy+\gyz\vz,\label{Equ:vel4y_d}\\
    u_z &= \gxz\vx+\gyz\vy+\gzz\vz,\label{Equ:vel4z_d}\\
    u_t &= -\alp W+\beta_x u_x+\beta_y u_y + \beta_z u_z.\label{Equ:vel4t_d}
  \end{align}
  The shift with the lowered components can be calculated via Equs.~\eqref{Equ:betx_d}-\eqref{Equ:betz_d}.
  \item Magnetic field components in fluid frame with upper indices: To obtain the magnetic field components in the fluid
  frame with upper indices $b^{\mu}$, we require the gauge variables $\alpha$ and $\beta^i$, as well as the undensitized
  magnetic field components in the Eulerian frame $B^i$, the Lorentz factor $W$, the 3-velocity $\tilde{u}^i$ and the spatial components of the 4-velocity $u^{\mu}$
  \begin{align}
    b^t&=\frac{B^iW\tilde{u}_i}{\alpha},\nonumber\\
    b^{i}&=\frac{B^i+\alpha b^tu^i}{W}.\nonumber
  \end{align}
  After computing $B^iW\tilde{u}_i$ via Equ.~\eqref{Equ:BWv}, undensitizing the Eulerian magnetic field components with
  Equs.~\eqref{Equ:undensmag1}-\eqref{Equ:undensmag3}, computing the Lorentz factor with Equ.~\eqref{Equ:lorentz}, and
  computing the four velocity components with Equs.~\eqref{Equ:vel4t}-\eqref{Equ:vel4z}, we can compute the
  magnetic field components in the fluid frame via
  \begin{align}
    b^t&=\frac{B^iW\tilde{u}_i}{\alp},\label{Equ:bt}\\
    b^x&=\frac{\Bx+\alp b^t u^x}{W},\label{Equ:bx}\\
    b^y&=\frac{\By+\alp b^t u^y}{W},\label{Equ:by}\\
    b^z&=\frac{\Bz+\alp b^t u^z}{W}.\label{Equ:bz}
  \end{align}
  \item Shift with lower components: Since AthenaK only outputs $\beta^i$, it is instructive for some computations to
  compute the shift with lower components, i.e. $\beta_i$, via
  \begin{align}
    \beta_i&=\gamma_{ik}\beta^k,\nonumber\\
    \beta_x&=\gamma_{xx}\beta^x+\gamma_{xy}\beta^y+\gamma_{xz}\beta^z,\nonumber\\
    \beta_y&=\gamma_{yx}\beta^x+\gamma_{yy}\beta^y+\gamma_{yz}\beta^z,\nonumber\\
    \beta_z&=\gamma_{xz}\beta^x+\gamma_{yz}\beta^y+\gamma_{zz}\beta^z.\nonumber
  \end{align}
  To achieve this, we require the metric components $\gamma_{ij}$ from the ouput, as well as the shift with
  upper components, or in terms of AthenaK variables we can write
  \begin{align}
    \beta_x &=\gxx\betx+\gxy\bety+\gxz\betz,\label{Equ:betx_d}\\
    \beta_y &=\gxy\betx+\gyy\bety+\gyz\betz,\label{Equ:bety_d}\\
    \beta_z &=\gxz\betx+\gyz\bety+\gzz\betz\label{Equ:betz_d}.
  \end{align}
  \item Amplitude of the shift vector: For the computation of the amplitude of the shift vector, we require the AthenaK
  shift with upper indices, i.e. $\beta^i$, as well as the metric components. It can then be
  computed via
  \begin{align}
    \beta^2 = \gamma_{ij}\beta^i\beta^j=\gamma_{xx}\beta^x\beta^x+2\gamma_{xy}\beta^x\beta^y+2\gamma_{xz}\beta^x\beta^z+
              \gamma_{yy}\beta^y\beta^y+2\gamma_{yz}\beta^y\beta^z+\gamma_{zz}\beta^z\beta^z.\nonumber
  \end{align}
  Or in terms of AthenaK variables
  \begin{align}\label{Equ:betamag}
    \beta^2=&\gxx\betx\betx+2\gxy\betx\bety+2\gxz\betx\betz\nonumber\\
    &+\gyy\bety\bety+2\gyz\bety\betz+\gzz\betz\betz.
  \end{align}
  \item Magnetic field components in fluid frame with lower indices: To obtain the magnetic field components in the fluid frame
  with lower indices $b_{\mu}$, we require the shift $\beta^i$, the magnetic field components in the fluid frame with upper inidices $b^{\mu}$,
  and the metric components $\gamma_{ij}$. The general lowering operation is given by
  \begin{align}
    b_t&=g_{\mu t}b^{\mu}=g_{tt}b^t+g_{tj}b^j=(-\alpha^2+\beta^2)b^t+\beta_jb^j,\nonumber\\
    b_j&=g_{\mu j}b^{\mu}=g_{tj}b^t+g_{ij}b^i=\beta_j b^t+\gamma_{ij}b^i.\nonumber
  \end{align}
  The latter equality is given by the line element in 3+1 notation
  \begin{align}
    ds^2&=-\alpha^2dt^2+\gamma_{ij}(dx^i+\beta^i dt)(dx^j+\beta^jdt),\nonumber \\
    g_{tt}&=-\alpha^+\beta_i\beta^i=-\alpha^2+\beta^2,\nonumber\\
    g_{ti}&=\beta_i,\nonumber\\
    g_{ij}&=\gamma_{ij}.\nonumber
  \end{align}
  After using Equs.~\eqref{Equ:bt}-\eqref{Equ:bz} to compute the magnetic field components in the fluid frame with
  upper indices, Equ.~\eqref{Equ:betamag} to compute the amplitude of the shift vector and Equs.~\eqref{Equ:betx_d}-\eqref{Equ:betz_d} to compute the shift vector with lower indices,
  we can use the AthenaK variables to get the lower indices of the magnetic field components in the fluid frame via
  \begin{align}
    b_t=&(-\alp^2+\beta^2)b^t+\beta_xb^x+\beta_yb^y+\beta_zb^z,\label{Equ:bt_d}\\
    b_x=&\gxx(b^t\betx+b^x)+\gxy(b^t\bety+b^y)+\gxz(b^t\betz+b^z),\label{Equ:bx_d}\\
    b_y=&\gxy(b^t\betx+b^x)+\gyy(b^t\bety+b^y)+\gyz(b^t\betz+b^z),\label{Equ:by_d}\\
    b_z=&\gxz(b^t\betx+b^x)+\gyz(b^t\bety+b^y)+\gzz(b^t\betz+b^z).\label{Equ:bz_d}
  \end{align}
  \item Local angular velocity: To compute the local angular velocity, we require the AthenaK 3-velocity $\tilde{u}^i$,
  the metric components $\gamma_{ij}$ and the gauge variables. In general we have
  \begin{align}\label{Equ:Omega}
    \Omega = \frac{u^{\phi}}{u^t}.
  \end{align}
  The velocity component $\phi$ is defined in terms of the Jacobian and the spatial components of the 4-velocity, i.e.
  \begin{align}
    u^{\phi}=\frac{\partial\phi}{\partial x^i}u^i.\nonumber
  \end{align}
  With $\phi=\arctan{\left(\frac{y}{x}\right)}$, we can compute the local angular velocity on any plane. For instance,
  in the xy-plane we find
  \begin{align}
    u^{\phi}&=\frac{\partial\phi}{\partial x}u^x+\frac{\partial \phi}{\partial y}u^y=-\frac{y}{x^2+y^2}u^x+\frac{x}{x^2+y^2}u^y=\frac{xu^y-yu^x}{R^2},\nonumber\\
    \Omega &= \frac{xu^y-yu^x}{R^2u^t}.\nonumber
  \end{align}
  Here, the 4-velocity components with upper indices can be computed via Equs.~\eqref{Equ:vel4t}-\eqref{Equ:vel4z}.
  \item Radial component of a 3-vector: For the computation of the radial component of a 3-vector with components $(p^x,p^y,p^z)$, we can write
  \begin{align}
    p^r = p^x\sin{\theta}\cos{\phi}+p^y\sin(\theta)\sin{\phi}+p^z\cos{\theta}.\label{Equ:radcomp}
  \end{align}
  \item Raise spatial vector: To raise a spatial vector with lower indices $(p_x,p_y,p_z)$, we can use the inverse spatial metric components from
  Equs.~\eqref{Equ:invgxx}-\eqref{Equ:invgyz} to get
  \begin{align}
    p^x &= \gamma&{xx}p^x+\gamma^{xy}p^y+\gamma^{xz}p^z,\label{Equ:raisepx}\\
    p^y &= \gamma&{xy}p^x+\gamma^{yy}p^y+\gamma^{yz}p^z,\label{Equ:raisepy}\\
    p^z &= \gamma&{xz}p^x+\gamma^{yz}p^y+\gamma^{zz}p^z.\label{Equ:raisepz}
  \end{align}
  \item Metric-normalized radial projection: To project a vector $p^i$ along $\frac{r_i}{\sqrt{g^{ij}r_ir_j}}$, we first compute the
  inverse spatial metric via Equs.~\eqref{Equ:invgxx}-\eqref{Equ:invgyz}. The inverse 4-metric $g^{ij}$ is then obtained via
  \begin{align}
    g^{xx} &= \gamma^{xx}-(\beta^x\beta^x)/\alpha^2,\nonumber\\
    g^{yy} &= \gamma^{yy}-(\beta^y\beta^y)/\alpha^2,\nonumber\\
    g^{zz} &= \gamma^{zz}-(\beta^z\beta^z)/\alpha^2,\nonumber\\
    g^{xy} &= \gamma^{xy}-(\beta^x\beta^y)/\alpha^2,\nonumber\\
    g^{xz} &= \gamma^{xz}-(\beta^x\beta^z)/\alpha^2,\nonumber\\
    g^{yz} &= \gamma^{yz}-(\beta^y\beta^z)/\alpha^2.\nonumber
  \end{align}
  The components of the vector $r_i$ are computed with the angles of the sphere
  \begin{align}
    r_x&=r\sin{\theta}\cos{\phi},\nonumber\\
    r_y&=r\sin{\theta}\sin{\phi},\nonumber\\
    r_z&=r\cos{\theta}\nonumber
  \end{align}
  The denominator $\sqrt{g^{ij}r_ir_j}$ is determined as
  \begin{align}
    \sqrt{g^{ij}r_ir_j}=\sqrt{g^{xx}r_x^2+g^{yy}r_y^2+g^{zz}r_z^2+2g^{xy}r_xr_y+2g^{xz}r_xr_z+2g_{yz}r_yr_z}.\nonumber
  \end{align}
  Lastly the projection of the vector $p^i$ is defined as
  \begin{align}
    p^r_{\mathrm{norm}}=\frac{r_xp^x+r_yp^y+r_zp^z}{\sqrt{g^{ij}r_ir_j}}.\label{Equ:radnormproj}
  \end{align}
\end{itemize}

\subsection*{Computations}
\step{Poynting flux}
For the computation of the Poynting flux across a shell, we require a set of output variables interpolated onto a sphere
of fixed radius and specified in the parameter file via:
\begin{verbatim}
<output1>
file_type = sph
variable  = velx
dt        = 10.0
radius    = 400.0
ntheta    = 256
\end{verbatim}
The variables we require are:
\begin{itemize}
  \item The ADM metric components \var{adm\_gxx}, \var{adm\_gxy}, \var{adm\_gxz},
  \var{adm\_gyy}, \var{adm\_gyz}, and \var{adm\_gzz} or the entire ADM output block \var{adm}. Note:
  if the entire \var{adm} block is specified, then the six extrinsic curvature components are also
  stored which are not needed for this computation and might fill unnecessary disk space.
  \item The gauge variables \var{z4c\_alpha}, \var{z4c\_betax}, \var{z4c\_betay}, and \var{z4c\_betaz}.
  \item The Lorentz-weighted 3-velocity $\tilde{u}^i$, given by the output variables \var{velx}, \var{vely}, and
  \var{velz}.
  \item The densitized magnetic field components in the Eulerian frame $\tilde{B}^i$, given by \var{bcc1}, \var{bcc2}, and
  \var{bcc3}.
\end{itemize}
A step-by-step approach is:
\begin{enumerate}
  \item Compute the Lorentz-factor via Equ.~\eqref{Equ:lorentz}.
  \item Compute the 4-velocity with upper indices $u^i$ using Equs.~\eqref{Equ:vel4t}-\eqref{Equ:vel4z}.
  \item Compute the lower time-component of the 4-velocity $u_t$ via Equ.~\eqref{Equ:vel4t_d}.
  \item Compute the determinant of the metric $\gamma$ via (Equ.~\eqref{Equ:detgamma}), to undensitize
  the magnetic field components with Equs.~\eqref{Equ:undensmag1}-\eqref{Equ:undensmag3}.
  \item Compute $B^2$ in the Eulerian frame via Equ.~\eqref{Equ:Bmag} and $B^iW\tilde{u}_i$ via Equ.~\eqref{Equ:BWv}.
  \item Compute $b^2$ in the fluid frame through Equ.~\eqref{Equ:bmag}.
  \item Compute the magnetic field components in the fluid frame $b^{\mu}$ via Equs.\eqref{Equ:bt_d}-\eqref{Equ:bz_d}.
  \item Lower the indices on the shift vector ($\beta_i$) via Equs.~\eqref{Equ:betx_d}-\eqref{Equ:betz_d}.
  \item Compute $\beta^2$ with Equ.~\eqref{Equ:betamag}.
  \item Compute the lower time-component of the magnetic field in the fluid frame, i.e. $b_t$, with Equ.~\eqref{Equ:bt_d}.
  \item Compute the radial component of the velocity vector $(u^x,u^y,u^z)$ and the magnetic field $(b^x,b^y,b^z)$ (in the fluid frame)
  with Equ.~\eqref{Equ:radcomp} ($u^r$ and $b^r$).
  \item Compute the integrand as
  \begin{align}
    I = (b^2u^ru_t-b^rb_t)\alpha\sqrt{\gamma}r^2\sin{\theta}.\nonumber
  \end{align}
  The spherical coordinates $(r,\theta,\phi)$ can be obtained from the spherical output \var{sph}.
  \item Lastly, integrate $I$ over the sphere with angles $(\theta,\phi)$ as dictated by the spherical output.
\end{enumerate}

\step{$\lambda_{\mathrm{MRI}}$ and quality factor $Q$}
For the computation of the $\lambda_{\mathrm{MRI}}$ parameter and the quality factor $Q$, we ideally require 3D
binary output. However, this procedure also works on a 2D slice. In general, we require
\begin{verbatim}
<output2>
file_type = bin
variable  = dens
dt        = 10.0
slice_x3  = 0.0 # optional
\end{verbatim}
The variables we require are:
\begin{itemize}
  \item The density $\rho$ given by \var{dens}, the temperature \var{temperature}, the electron fraction $Y_e$ \var{s\_00},
  the densitized magnetic field components \var{bcc1}, \var{bcc2}, \var{bcc3}, and the Lorentz-weighted 3-velocity components $\tilde{u}^i$, given
  by \var{velx}, \var{vely}, and \var{velz}.
  \item The ADM metric components \var{adm\_gxx}, \var{adm\_gxy}, \var{adm\_gxz},
  \var{adm\_gyy}, \var{adm\_gyz}, and \var{adm\_gzz} or the entire ADM output block \var{adm}. Note:
  if the entire \var{adm} block is specified, then the six extrinsic curvature components are also
  stored which are not needed for this computation and might fill unnecessary disk space.
  \item The gauge variables \var{z4c\_alpha}, \var{z4c\_betax}, \var{z4c\_betay}, and \var{z4c\_betaz}.
\end{itemize}
A step-by-step approach is:
\begin{enumerate}
  \item Compute the Lorentz-factor via Equ.~\eqref{Equ:lorentz}.
  \item Compute the determinant of the metric $\gamma$ via (Equ.~\eqref{Equ:detgamma}), to undensitize
  the magnetic field components with Equs.~\eqref{Equ:undensmag1}-\eqref{Equ:undensmag3}.
  \item Compute $B^2$ in the Eulerian frame via Equ.~\eqref{Equ:Bmag} and $B^iW\tilde{u}_i$ via Equ.~\eqref{Equ:BWv}.
  \item Compute $b^2$ in the fluid frame through Equ.~\eqref{Equ:bmag}.
  \item Compute the magnetic field components in the fluid frame $b^{\mu}$ via Equs.\eqref{Equ:bt_d}-\eqref{Equ:bz_d}.
  \item Compute the 4-velocity with upper indices $u^i$ using Equs.~\eqref{Equ:vel4t}-\eqref{Equ:vel4z}.
  \item Compute the spatial magnetic field components in the fluid frame with lower indices $b_i$ via
  Equs.~\eqref{Equ:bx_d}-\eqref{Equ:bz_d}.
  \item Compute the local angular velocity either on the plane or if 3D over the entire domain via Equ.~\eqref{Equ:Omega}.
  \item The components of $\lambda_{\mathrm{MRI}}$ are then computed via
  \begin{align}
    \lambda_{\mathrm{MRI},i}=\frac{2\pi b_i}{\Omega\sqrt{\rho h+b^2}},\nonumber
  \end{align}
  where $h$ is the enthalpy given by the equation of state and the state vector $(n_B,Y_e,T)$ at a given point.
  \item The quality average quality factor for each direction is then
  \begin{align}
    Q_i = \langle\lambda_{\mathrm{MRI},i}/\Delta x\rangle,\nonumber
  \end{align}
  where $\Delta x$ is the grid spacing at the given grid point.
\end{enumerate}

\step{Ejecta - Bernoulli and Geodesic criterion}
To compute the ejecta properties based on the Bernoulli and Geodesic criterion, we need data interpolated onto a sphere
at a fixed radius, i.e.
\begin{verbatim}
<output3>
file_type = sph
variable  = velx
dt        = 10.0
radius    = 400.0
ntheta    = 256
\end{verbatim}
The variables we require are:
\begin{itemize}
  \item The density $\rho$ given by \var{dens}, the temperature \var{temperature}, the electron fraction $Y_e$ \var{s\_00},
  and the Lorentz-weighted 3-velocity components $\tilde{u}^i$, given by \var{velx}, \var{vely}, and \var{velz}.
  \item The ADM metric components \var{adm\_gxx}, \var{adm\_gxy}, \var{adm\_gxz},
  \var{adm\_gyy}, \var{adm\_gyz}, and \var{adm\_gzz} or the entire ADM output block \var{adm}. Note:
  if the entire \var{adm} block is specified, then the six extrinsic curvature components are also
  stored which are not needed for this computation and might fill unnecessary disk space.
  \item The gauge variables \var{z4c\_alpha}, \var{z4c\_betax}, \var{z4c\_betay}, and \var{z4c\_betaz}.
\end{itemize}
We also need the 3D equation of state with columns $(n_B,Y_e,T)$. This table will be used to calculate the specific
enthalpy. A step-by-step approach is:
\begin{enumerate}
  \item Take the points for the density, electron fraction, and the temperature from the \var{sph} output and clip the points
  to the ranges of the 3D table. Change the density and temperature columns to log-scale.
  \item Set up interpolators to interpolate the total energy density $\frac{e}{n_Bm_n}$ and the pressure $\frac{p}{n_Bm_n}$ from the table given $(\log_{10}(n_B),Y_e,\log_{10}(T))$.
  \item For all the points on the spherical shell, interpolate to obtain the total energy density and the pressure and then compute the
  specific enthalpy via
  \begin{align}
    h = 1+e_{\mathrm{intp}}+p_{\mathrm{intp}}.\nonumber
  \end{align}
  \item Compute the Lorentz-factor via Equ.~\eqref{Equ:lorentz}.
  \item Compute the 4-velocity with upper indices $u^i$ via Equs.~\eqref{Equ:vel4t}-\eqref{Equ:vel4z}.
  \item Compute the radial component from the spatial part of the 4-velocity with Equ.~\eqref{Equ:radcomp}.
  \item Lower the index on the time-component of the 4-velocity $u_t$ via Equ.~\eqref{Equ:vel4t_d}.
  \item Create masks, such that the density is above the density floor and that the radial component of the velocity
  $u^r$ is greater than 0.
  \item Compute a criterion mask with
  \begin{align}
    \mathrm{Bernoulli:}&\quad h\cdot u_t < h_{\mathrm{min}},\nonumber\\
    \mathrm{Geodesic:}&\quad u_t < -1.\nonumber
  \end{align}
  \item Compute the determinant of the spatial metric $\gamma$ with Equ.~\eqref{Equ:detgamma} and the base
  integrand via
  \begin{align}
    I = \rho u^r\sqrt{\gamma}r^2\sin{\theta}.\nonumber
  \end{align}
  \item See the documentation on the ejecta and neutrino analysis in the repo for more information about the various quantities one
  can calculate with the above steps.
\end{enumerate}

\step{Neutrino luminosity and average energy density}
For the neutrino luminosity and average energy density we again require spherical output at a fixed radius, i.e.
\begin{verbatim}
<output4>
file_type = sph
variable  = rad_m1_F
dt        = 10.0
radius    = 400.0
ntheta    = 256
\end{verbatim}
The variables we need are:
\begin{itemize}
  \item The ADM metric components \var{adm\_gxx}, \var{adm\_gxy}, \var{adm\_gxz},
  \var{adm\_gyy}, \var{adm\_gyz}, and \var{adm\_gzz} or the entire ADM output block \var{adm}. Note:
  if the entire \var{adm} block is specified, then the six extrinsic curvature components are also
  stored which are not needed for this computation and might fill unnecessary disk space.
  \item The gauge variables \var{z4c\_alpha}, \var{z4c\_betax}, \var{z4c\_betay}, and \var{z4c\_betaz}.
  \item The neutrino energy density per species \var{rad\_m1\_E}, the neutrino number flux per species \var{rad\_m1\_F}, and
  the neutrino number density per species \var{rad\_m1\_N} in the Eulerian frame.
\end{itemize}
A step-by-step guide is:
\begin{enumerate}
  \item Take the neutrino number flux $F_i$ and raise its components with the inverse spatial metric, which can
  be computed with Equs.~\eqref{Equ:invgxx}-\eqref{Equ:invgyz}. The raising operation is then defined via Equs.~\eqref{Equ:raisepx}-\eqref{Equ:raisepy}.
  We obtain the vector $F^i$.
  \item Taking the raised neutrino number flux vector in the Eulerian frame $F^i$, we can compute the metric-normalized radial projection $F^r_{\mathrm{norm}}$ via
  Equ.~\eqref{Equ:radnormproj}.
  \item The same projection is done for the shift vector $\beta^i$ to obtain $\beta^r_{\mathrm{norm}}$.
  \item The radial flux is compute as
  \begin{align}
    F_{\mathrm{radial}}=\alpha F^r-\beta^rE.\nonumber
  \end{align}
  \item Compute the neutrino luminosity per species $s$ as
  \begin{align}
    L^{\mathrm{code}}_{s}=\sum F_{\mathrm{radial}}w,\nonumber
  \end{align}
  where $w$ are the surface weights of the extraction sphere. The luminosity has to be converted into physical units afterwards.
  \item Compute the average energy density of the neutrino species, by computing the number flux as
  \begin{align}
    N_{\mathrm{flux}}=\frac{F_{\mathrm{radial}}N_{\mathrm{code}}}{E},\nonumber
  \end{align}
  where $N_{\mathrm{code}}$ is the neutrino number density in code units.
  \item Lastly, compute the change of the neutrino number density as
  \begin{align}
    \Dot{N}_{\mathrm{code}}=\sum N_{\mathrm{flux}}w,\nonumber
  \end{align}
  and then compute the average energy density in code units as
  \begin{align}
    \langle\varepsilon\rangle_s=\frac{L^{\mathrm{code}}_s}{\Dot{N}_{\mathrm{code}}}.\nonumber
  \end{align}
\end{enumerate}
Note: the exact M1 number flux uses $f_{\nu}^a$, so the above is an approximation.

\step{M1 variables in fluid frame}
Since all the M1 variables in AthenaK are given in the Eulerian frame, we need to perform a set of
projections to obtain them in the fluid frame. This is particularly desireable if we want to do a
tracer analysis in post-processing. To convert the M1 variables, namely the radiation energy density $E$, the radiation flux $F^{\alpha}$,
and the radiation pressure tensor $P^{\alpha\beta}$ in the Eulerian frame into the fluid frame, we first
note that the definition of the radiation stress-energy tensor in both frames is
\begin{align}
  \mathrm{Eulerian\ frame:}&\quad T^{\alpha\beta}=En^{\alpha}n^{\beta}+F^{\alpha}n^{\beta}+n^{\alpha}F^{\beta}+P^{\alpha\beta},\\
  \mathrm{Fluid\ frame:}&\quad T^{\alpha\beta}=Ju^{\alpha}u^{\beta}+H^{\alpha}u^{\beta}+u^{\alpha}H^{\beta}+K^{\alpha\beta},
\end{align}
where $n^{\alpha}$ is the unit-normal to the t = const. hypersurface, $u^{\alpha}$ is the fluid four-velocity,
$J$ is the radiation energy density, $H^{\alpha}$ the radiation flux, and $K^{\alpha\beta}$ the radiation pressure
tensor in the fluid frame. We will need to carry out a set of projections to obtain the fluid frame variables. An additional
quantity in the fluid frame will be the comoving number density $n=\frac{N}{\Gamma}$, where $N$ is the number density
in the Eulerian frame and $\Gamma$ the ratio of neutrinos in the Eulerian and the fluid frame. Since this is not an evolved
variable, we can simply derive it on the way from the evolved ones. Note: the procedure has to be done per species.
A step-by-step approach is:
\begin{enumerate}
  \item Compute the lower and upper spacetime metric from the spatial metric $\gamma_{ij}$ and the gauge variables
  $\alpha$, and $\beta^i$ via
  \begin{align}
    g_{00}&=-\alpha^2+\beta^i\beta_i,\nonumber\\
    g_{0j}&=\beta_j,\nonumber\\
    g_{ij}&=\gamma_{ij},\nonumber\\
    g^{00}&=-\frac{1}{\alpha^2},\nonumber\\
    g^{0j}&=\frac{\beta^{j}}{\alpha^2},\nonumber\\
    g^{ij}&=\gamma^{ij}-\frac{\beta^i\beta^i}{\alpha^2}.\nonumber
  \end{align}
  These can be computed via the functions \var{adm::SpacetimeMetric} and \var{adm::SpacetimeUpperMetric} in \var{adm.hpp}.
  \item Populate a four-vector $n_{\alpha}$ using \var{pack\_n\_d} inside of \var{radiation\_m1\_tensors.hpp}, which defines
  \begin{align}
    n_0&=-\alpha,\nonumber\\
    n_i&=0.
  \end{align}
  \item Populate a four-vector $\beta^{\mu}$ with the spatial components of the shift $\beta^i$, where for $\mu=0$ the vector is
  zero (via the function \var{pack\_beta\_u} in \var{radiation\_m1\_tensors.hpp}).
  \item Compute the Lorentz-factor using
  Equ.~\eqref{Equ:lorentz}, and the Lorentz-weighted 3-velocity in AthenaK by using the function \var{get\_w\_lorentz} inside
  \var{radiation\_m1\_tensors.hpp}.
  \item Pack a vector for the 4-velocity $u^{\alpha}$ via \var{pack\_u\_u} with
  \begin{align}
    u^{0}&=\frac{W}{\alpha},\nonumber\\
    u^{i}&=\tilde{u}^i-\frac{W\beta^i}{\alpha},\nonumber
  \end{align}
  and pack a vector for the 3-velocity $v^{\alpha}$ (only spatial components important) via \var{pack\_v\_u} with
  \begin{align}
    v^{0}&=0,\nonumber\\
    v^{i}&=\frac{u^{i}}{W}+\frac{\beta^{i}}{\alpha}.\nonumber
  \end{align}
  Above, $\tilde{u}^i$ is the Lorentz-weighted 3-velocity that AthenaK uses.
  \item Use the \var{tensor\_contract} function in \var{radiation\_m1\_tensors.hpp} to compute $u_{\alpha}$ and
  $v_{\alpha}$ using the lower spacetime metric $g_{\alpha\beta}$ from step 1.
  \item Calculate the fluid projector $P^{\alpha}_{\beta}$ from $u_{\alpha}$ and $u^{\alpha}$ using \var{calc\_proj} inside \var{radiation\_m1\_helper.hpp} which
  computes
  \begin{align}
    P^{\alpha}_{\beta}=\delta^{\alpha}_{\beta}+u^{\alpha}u_{\beta},\nonumber
  \end{align}
  where $\delta^{\alpha}_{\beta}$ is the Kronecker-delta.
  \item For each neutrino species, get the radiation energy density $E$ from the array holding the solution and
  pack the radiation flux $F^{\alpha}$ using \var{pack\_F\_d} in \var{radiation\_m1\_tensors.hpp}, which computes
  \begin{align}
    F_{0}&=\beta^iF_i,\nonumber\\
    F_{i}&=F_i,\nonumber
  \end{align}
  where $F_{i}$ is the radiation flux per species evolved on the given hypersurface.
  \item Using the spacetime metric with lower indices $g_{\alpha\beta}$, with upper indices $g^{\alpha\beta}$,
  the unit-normal with lower index $n_{\alpha}$, the Lorentz-factor $W$, the fluid 4-velocity $u^{\alpha}$, the fluid 3-velocity
  $v_i$, the fluid projector $P^{\alpha}_{\beta}$, the radiation energy density $E$, and the radiation flux $F_{\alpha}$ in the Eulerian frame
  plus the Eddington-factor $\chi$, we can compute the radiation stress-energy tensor $P_{\alpha\beta}$ using the \var{apply\_closure}
  function inside \var{radiation\_m1\_helpers.hpp}. This function computes
  \begin{align}
    P_{\alpha\beta}^{\mathrm{thin}}=&\frac{E}{F_{\mu}F^{\mu}}F_{\alpha}F_{\beta},\nonumber\\
    \gamma^{\kappa}_{\alpha}H_{\kappa}=&\frac{F_{\alpha}}{W}+\frac{W}{2W^2+1}v_{\alpha}((4W^2+1)F^{\mu}v_{\mu}-4W^2E),\nonumber\\
    P_{\alpha\beta}^{\mathrm{thick}}=&\left(\frac{2W^2-1}{2W^2+1}E-2W^2F^{\mu}v_{\mu}\right)(4W^2v_{\alpha}v_{\beta}+g_{\alpha\beta}+n_{\alpha}n_{\beta})\nonumber\\
    &+W(\gamma^{\kappa}_{\alpha}H_{\kappa}v_{\beta}+\gamma^{\kappa}_{\beta}H_{\kappa}v_{\alpha}),\nonumber\\
    P_{\alpha\beta}=&\frac{3\chi-1}{2}P_{\alpha\beta}^{\mathrm{thin}}+\frac{3(1-\chi)}{2}P_{\alpha\beta}^{\mathrm{thick}}.\nonumber
  \end{align}
  \item Using the radiation pressure tensor in the Eulerian frame $P_{\alpha\beta}$, we can now compute the stress-energy tensor
  $T_{\alpha\beta}$ in the Eulerian frame using the function \var{assemble\_rT} in \var{radiation\_m1\_helpers.hpp}, which computes
  \begin{align}
    T_{\alpha\beta}=En_{\alpha}n_{\beta}+F_{\alpha}n_{\beta}+n_{\alpha}F_{\beta}+P_{\alpha\beta}.\nonumber
  \end{align}
  \item With the stress-energy tensor in the Eulerian frame $T_{\alpha\beta}$, we can now contract this with the four-velocity vector
  to obtain the radiation energy density in the fluid frame via
  \begin{align}
    \boxed{J=T_{\alpha\beta}u^{\alpha}u^{\beta}}.
  \end{align}
  This computation is encapsulated inside the function \var{calc\_J\_from\_rT} inside \var{radiation\_m1\_helpers.hpp}.
  \item Using the fluid projector $P^{\alpha}_{\beta}$, the stress-energy tensor in the Eulerian frame $T_{\alpha\beta}$ and
  the fluid four-velocity $u^{\alpha}$, we can compute the radiation flux in the fluid frame $H_{\alpha}$ via
  \begin{align}
    \boxed{H_{\alpha}=-P^{\beta}_{\alpha}u^{\kappa}T_{\beta\kappa}}.
  \end{align}
  This computation is encapsulated inside the function \var{calc\_H\_from\_rT} inside \var{radiation\_m1\_helpers.hpp}.
  Contract with the spacetime metric $g^{\alpha\beta}$ to obtain $H^{\alpha}$.
  \item Enforce that $J$ and $H_{\alpha}$ are above their respective floors with \var{apply\_floor} inside \var{radiation\_m1\_helpers.hpp}.
  \item Compute the comoving number density $n$, using $J$, $E$, and $F_{\alpha}$. For that we compute the ratio of neutrinos in the Eulerian and
  the fluid frame $\Gamma$ via \var{compute\_Gamma} inside {radiation\_m1\_helpers.hpp} which defines
  \begin{align}
    \Gamma=W\frac{E}{J}(1-\frac{F_{\alpha}v^{\alpha}}{E}),\quad \boxed{n=\frac{N}{\Gamma}},
  \end{align}
  where $N$ is the number density in the Eulerian frame, and $v^{\alpha}$ the 3-velocity packed into a 4-vector.
  \item Optional: compute the unit-norm radiation number current $f_{\nu}$ using the fluid four-velocity $u^{\alpha}$,
  the radiation energy density $J$, and the radiation flux in the fluid frame $H^{\alpha}$ with
  \begin{align}
    \boxed{f_{\nu}^{\alpha}=u^{\alpha}+\frac{H^{\alpha}}{J}}.
  \end{align}
  The computation is inside the function \var{assemble\_fnu} in \var{radiation\_m1\_helpers.hpp}. Note: inside the computation there is
  a check that the radiation energy density in the fluid frame $J$ is above the floor for the radiation energy density in the Eulerian frame $E$. If this is
  not the case, the second term vanishes.
\end{enumerate}
\end{document}
